\section{Morale}

Morale represents how well-trained and disciplined a detachment is. It also reflects the detachment's ability to withstand casualties before its discipline breaks and it falls back.

\subsection{Affected Detachment}

A formation that loses at least half of its current bases or Wounds as the result of a single action, either an assault or the resolution of a shooting attack, becomes \textit{Affected}. \textit{(Inc)}

A detachment that becomes Affected must immediately make a Morale Test.

\subsection{Broken Detachment}

A detachment that has lost half or more of its initial Wounds is considered \textit{Broken}.

A Broken formation awards its Destruction Point value to the opponent. These points are used in certain scenarios. \textit{(Inc)}

\subsection{Morale Test}

Certain circumstances and abilities require a detachment to make a Morale Test. A detachment may be required to make several Morale Tests during the same turn if the appropriate circumstances arise.

Regardless of what caused the test, the procedure is always the same. Roll 2D6 and add 1 if the detachment has not lost any bases or Wounds. Additional bonuses or penalties may apply when resolving an assault.

If the result is equal to or greater than the detachment's Morale value, the test is passed.

If the test is failed, replace the detachment's current order with a Fall Back order, and it must immediately make a Fall Back move. If the detachment has not yet acted during the current turn, it loses the opportunity to do so.

A detachment without a Morale value automatically passes all Morale Tests.

\subsection{Falling Back}

A detachment that fails a Morale Test is considered to be \textit{Falling Back}.

A detachment with a Fall Back order cannot perform any actions during the Combat Phase. It may remove its Fall Back order by passing a Rally Test or during its next activation in the Movement Phase by performing one of the actions permitted by the Fall Back order.

If a Falling Back detachment is completely pinned and cannot be activated during the Movement Phase, it remains Falling Back for the rest of the turn.

While a detachment has a Fall Back order:

\begin{itemize}
    \item It has no zone of control.
    \item It cannot capture objectives.
    \item It cannot make Consolidation moves.
    \item Every base in the detachment suffers a $-2$ penalty to its AF.
\end{itemize}

\subsection{Rally Test}

When a detachment makes a Rally Test, it performs a Morale Test.

If the test is passed, remove its Fall Back order counter. The detachment may receive orders normally during the following turn.

If the test is failed, it retains its Fall Back order.

\subsection{Fall Back Movement}

A detachment making a Fall Back move moves between 0 and 25~cm, regardless of its Movement characteristic, subject to the following restrictions:

\begin{itemize}
    \item The detachment must leave every enemy interdiction zone by the shortest possible route. If this is impossible, it must move as far away as possible from all enemy bases.

    \item Once the detachment has left an enemy interdiction zone, it cannot enter another enemy interdiction zone.

    \item The detachment cannot enter an enemy zone of control, even if the enemy base is of a lower class.
\end{itemize}

Apart from these restrictions, the detachment may use its movement freely, including remaining in place.

A detachment engaged in an assault never makes a Fall Back move, whether or not it is pinned.

If, for any reason, the affected bases cannot leave enemy zones of control, they are destroyed.